\section{Síntesis y ampliación}

\subsection{Puntos clave}
\begin{itemize}
    \item «Dive into Kimi K2.5» enmarca toda la clase en tres líneas principales: data, algorithm e infra, lo que facilita desglosar el informe técnico de manera lógica.
    \item Mezclar datos visuales desde una etapa temprana y usar un pipeline sintético impulsado por benchmarks es la base que le permite a K2.5 conectar sus capacidades en imagen y texto, código y GUI.
    \item El GRM se convierte en la recompensa unificada; las tres etapas de entrenamiento (Preentrenamiento/SFT/RL) se superponen y generan una capacidad estable, y el patrón de deg-and-recover es una trayectoria de convergencia normal.
    \item El atractivo del agent swarm está en el Allcastrater aprendible, los subagents paralelos y las restricciones de recursos de los critical steps, lo que hace que múltiples Agents sean más eficientes que un solo Agent.
    \item La Infrastructure se evalúa mediante Critical Steps + Kimi code bench; aunque el docente explicó poco al respecto, las curvas de entrenamiento y el bandmark ya demuestran que el sistema es viable.
    \item La parte práctica muestra cómo combinar subtítulos, fotogramas clave y slides en un documento con una summary table y una jerarquía de boxes, para garantizar que pase el audit.
    \item Al final, la multimodalidad nativa más el Agent Swarm le dan a K2.5 la capacidad de ejecución de una «Visual Agency» real, y no solo comprensión visual.
\end{itemize}

\subsection{Tabla de síntesis}
\begin{table}[H]
\footnotesize
\centering
\begin{tabular}{p{3cm}p{7cm}p{4cm}}
\toprule
Dimensión & Idea clave & Evidencia \\
\midrule
Data & Mezclar datos visuales desde el inicio, definir benchmarks y luego usar un pipeline sintético para expandir la capacidad por etapas. & El experimento de proporciones 1:9 a 1:1 en la Figure 9 del apéndice; el proceso de construcción de datos en tres pasos que enumeró el docente. \\
Algorithm & GRM + PARL RL unifican la recompensa; el entrenamiento multietapa le da estabilidad a las tareas open-ended; Deg \& Recover es una fluctuación natural. & La explicación en el knowledgebox y el warningbox del GRM, el PARL agent loop y las curvas de entrenamiento de los critical steps. \\
Agent \& Infra & El Allcastrater garantiza la asignación de subagents, los Critical Steps limitan los recursos, y la Evaluation usa Kimi code bench y bandmark. & La explicación de createSubagents/assignTask en la Agent section, más la mención del bandmark de evaluación en la infra section. \\
Practice & El capítulo práctico estandariza el procesamiento del material, la inserción de fotogramas clave y slides, y la summary table, para que el documento cumpla con los requisitos del audit. & Los knowledgebox e importantbox del capítulo práctico, además de los ejemplos de figure timecode. \\
\bottomrule
\end{tabular}
\caption{La industrialización por capas, de los datos y el algoritmo hasta el agent/infra, refleja el cierre lógico de la arquitectura de multimodalidad nativa + Agent Swarm de K2.5.}
\end{table}

\subsection{Lecturas adicionales}
\begin{itemize}
    \item \href{https://www.kimi.com/ai-models/kimi-k2-5/}{Kimi K2.5 | Open Visual Agentic Model for Real Work} — Resumen oficial del modelo y descripción de la interfaz.
    \item \href{https://www.kimi.com/blog/kimi-k2-5}{Kimi K2.5 Tech Blog: Visual Agentic Intelligence} — Presenta en detalle la Visual Agency Intelligence y la estrategia de datos.
    \item \href{https://artificialanalysis.ai/articles/kimi-k2-5-everything-you-need-to-know}{Kimi K2.5 — Everything you need to know (Artificial Analysis)} — Un resumen externo de los benchmarks y el agent swarm.
    \item \href{https://futuretechdiaries.com/article.php?slug=kimi-k2-5-a-deep-dive-into-chinas-long-context-ai-model}{Kimi K2.5: A Deep Dive into China’s Long-Context AI Model} — Comparación y análisis de casos de uso.
    \item \href{https://platform.kimi.ai/docs/guide}{Kimi API Platform Guide} — Explica el método práctico para invocar agents visuales y usar la cadena de herramientas.
\end{itemize}

\subsection{Resumen de la sección}
Esta clase integradora finalmente nos muestra que K2.5 construye un agentic stack nativamente multimodal, en el que los datos, el algoritmo impulsado por el GRM y la coordinación de infra/agent forman un ciclo cerrado que le da al modelo una auténtica Visual Agency Intelligence.
%% --- Fin del contenido --- %%

\end{document}
